\chapter{Current Status}

\section{Implementation phases}

The architecture described in \texttt{charlotte-sitas-xous-architecture.md} is
substantially implemented and boot-validated on QEMU (\texttt{-M
virt,gic-version=3}, Apple Silicon HVF or TCG).

\begin{longtable}{|c|p{9cm}|c|}
\hline
\textbf{Phase} & \textbf{Description} & \textbf{Status} \\
\hline
1 & Capability table, rights masks, object teardown & Done \\
2 & Endpoint IPC v1: scalar send/call/receive/reply, connection delegation,
    reply tokens & Done \\
3 & Userspace name/service manager, bootstrap delivery, ELF loader, supervisor,
    generation tracking, long names via memory objects & Done \\
4 & First-class memory objects: allocate, map, unmap, close, ownership
    accounting & Done \\
5 & Memory IPC: Copy, Move, BorrowRead, BorrowWrite, reply-bound revocation,
    cancellation, server-death recovery & Done \\
6 & CQ sub-system normalisation: operation IDs, detached submission,
    CQ\_WAIT/CQ\_WAKE, per-shard CQ partitioning, backlog batching, \S8.2
    32-byte richer completion records & Done \\
7 & Sitas endpoint/CQ backend: endpoint readiness binding, unified shard
    wait, ShardExecutor (budgeted polling, task wakeup from drained events),
    ShardParker seam (spin free), per-shard CQ rings, kv::spin\_recv
    retired & Done \\
8 & Userspace UART driver: delegated MMIO + IRQ, EL0 MMIO writes,
    interrupt-driven deferred reads, driver crash $\rightarrow$ device reset
    $\rightarrow$ outstanding-op reconciliation $\rightarrow$ generation-2
    restart & Done \\
9 & Virtio-net driver: PCI discovery, BAR0 + IRQ delegation, virtio init
    sequence, MAC read, virtqueue setup, frame TX/RX (compiles),
    MEMORY\_GET\_PHYS syscall. Protocol crates extracted. Smoltcp 0.13
    adapter + TCP/IP service binary (compile). Runtime validation blocked by
    HVF EL0-MMIO limitation. & Driver built, \\
    & & pending validation \\
10 & Lock-free device interrupt delivery (deferred wake), idempotent scheduler
    wakes & Done \\
\hline
\end{longtable}

\section{Success criteria}

\begin{itemize}
    \item[1.] Two EL0 domains communicate through a bounded endpoint
        without ASID/LP.
    \item[2.] Deferred async call while a shard continues other tasks.
    \item[3.] Moved memory object --- sender locked out until return.
    \item[4.] Lending enforced by mappings, including death cleanup.
    \item[5.] Single \texttt{CQ\_WAIT} for endpoint + completions +
        device interrupts.
    \item[6.] Terminal CQ results survive ring overflow (non-lossy
        backlog).
    \item[7.] Coalesced remote shard wakes (idempotent scheduler fix).
    \item[8.] Userspace UART driver with only delegated MMIO + IRQ.
    \item[9.] Restart invalidates stale connections + device reset +
        op reconciliation.
    \item[10.] \textasciitilde{} Virtio-net data path with batched packet
        buffers (driver compiles, runtime validation pending).
    \item[11.] \textasciitilde{} Fixed-width stable ABI representations (CQ
        entries are 32-byte \S8.2; not yet formally audited).
    \item[12.] \textasciitilde{} Capability misuse and cancellation races have
        adversarial tests (six IPC error-path tests added; not yet
        comprehensive).
\end{itemize}

\section{Verified (boots on QEMU, zero panics)}

\begin{itemize}
    \item \textbf{AArch64:} boots, all self-tests pass, the full EL0 service
        suite (name service, echo, client, UART driver, sitas executor) runs
        through \texttt{crt0/cmain}. The sitas \texttt{basic\_kv} binary runs
        with 2 shards on per-shard CQ rings and reports \texttt{total\_len == 3}.
    \item \textbf{x86-64:} boots with \texttt{-cpu max}, core self-tests pass.
        (x86-64 ring-3 userspace is implemented but not yet exercised with a
        user thread.)
    \item \textbf{Completion sub-system:} submit, complete, poll, wait, cancel,
        close, \texttt{submit\_worker} (exit-observer), CQ ring, CQ ring
        integrated with AS, non-lossy CQ overflow backlog, per-shard CQ
        partitioning, \S8.2 32-byte entries with operation/cookie/status/result.
    \item \textbf{Endpoint IPC:} bounded server queues, scalar messages,
        connection delegation with rights attenuation, memory-object attachments
        (move/lend/copy), deferred replies, reply-time connection delegation,
        combined connection+memory attachments.
    \item \textbf{Device capabilities:} MmioRegion and Interrupt object types,
        user-accessible Device-nGnRnE page mapping, GICv3 SPI enable/route/mask,
        IRQ$\rightarrow$CQ coalesced delivery (lock-free deferred wake).
    \item \textbf{Userspace drivers:} reference UART driver (PL011) with
        delegated MMIO + IRQ, interrupt-driven deferred reads, crash-test
        restart with device reset and outstanding-operation reconciliation.
    \item \textbf{Syscall dispatch:} AArch64 SVC handler dispatches 49 syscalls
        across completion, IPC, memory, device, and CQ operations. Caller ASID
        is derived from the trap frame, never a userspace parameter.
    \item \textbf{Scheduler:} RoundRobin, quantum timer, block/yield/wake,
        idempotent wakes (unified shard wait), deferred reaper, per-LP thread
        pinning (\texttt{submit\_to\_lp}).
    \item \textbf{sitas integration:} sitas-core (no\_std executor) compiles for
        aarch64-unknown-none. sitas-charlotte (ReactorBackend + ShardRuntime)
        links against the kernel's syscall ABI. Per-shard CQ rings are mapped
        by the loader contract. ShardParker seam (CQ\_WAIT/CQ\_WAKE) retires
        busy-waiting; ShardExecutor implements budgeted polling with task
        wakeup from drained events.
    \item \textbf{Protocol crates:} \texttt{charlotte-protocol-net}
        (frame-level NIC driver protocol), \texttt{charlotte-protocol-msg}
        (reliable message layer with sequenced/acked delivery).
    \item \textbf{Networking:} smoltcp 0.13 adapter (implements
        \texttt{phy::Device} over the NIC-driver endpoint) and TCP/IP service
        binary compile for aarch64-none.
    \item \textbf{CI:} GitHub Actions builds both architectures with
        \texttt{-D warnings}, QEMU boot gate.
\end{itemize}

\section{Repositories and branches}

\begin{itemize}
    \item \textbf{CharlotteOS} (\texttt{dev}): consolidated kernel ---
        completion sub-system, endpoint IPC, memory objects, device
        capabilities, syscall dispatch, bounded IPI, ShardLocal, ShardMailbox,
        CQ ring, EL0 service loader, supervisor, self-test suite (adversarial
        IPC, device-capability, UART driver), developer manual.
    \item \textbf{sitas} (\texttt{reactor-handle-seam}): \texttt{ReactorBackend}
        trait, \texttt{ShardExecutor} (futures executor with budgeted polling),
        waker-integrated channels, \texttt{ShardParker} seam, per-shard CQ
        reactor/waker/parker targeting, \texttt{sitas-core} no\_std crate,
        \texttt{sitas-charlotte} crate (ReactorBackend + ShardRuntime backend).
\end{itemize}

\section{Known limitations}

\begin{itemize}
    \item \textbf{x86-64 ring-3 userspace:} SYSCALL entry trampoline exists; the
        EL0 test and real-EL0 user thread are AArch64-only.
    \item \textbf{Virtio-net runtime:} the driver and TCP/IP stack compile but
        cannot be runtime-validated on macOS (HVF does not emulate EL0 MMIO
        instructions). Validation requires a Linux KVM host.
    \item \textbf{General process loader:} the ELF loader is fixed-address for
        the reference service images, not a general loader with
        argv/env/auxv-style setup.
    \item \textbf{Resource accounting:} per-domain quotas are not yet
        implemented; capability-table capacity is fixed.
\end{itemize}

\section{Build and boot}

\textbf{AArch64 (macOS, HVF):}
\begin{verbatim}
cargo build --package catten \
  --target target_specs/aarch64-unknown-none-catten.json
./scripts/qemu-boot-macos.sh 60 2
\end{verbatim}

\textbf{sitas user binary (requires nightly + build-std):}
\begin{verbatim}
cd crates/catten-user && \
cargo +nightly build --target aarch64-unknown-none.json \
  -Zbuild-std=core,alloc
\end{verbatim}

\textbf{EL0 service programs:}
\begin{verbatim}
cd crates/catten-services && \
cargo build --release --target aarch64-unknown-none.json \
  -Zbuild-std=core,alloc
\end{verbatim}

\endinput
